\printedpage{54}
% head: A. Grothendieck : Dérivateurs, Ch. XVIII, Édition des Universités de Montpellier II et Paris VII

\textsc{Proof} of 4.7.1. The implications a $\Longrightarrow$ a$'$ $\Longrightarrow$ a$''$ are tautological, as is b $\Longrightarrow$ b$'$. The implication b$'$ $\Longrightarrow$ c$'$ comes from the stability of $C^{\wedge}_{\pi}$ under $\varinjlim$ of cardinality $\leqslant \pi$, and from the fact that the $\varinjlim$ of the type under consideration (in particular, large relative to $\pi_0$) \textit{in} $\mathcal{M}$ are the same as in $C^{\wedge}$, i.e. are computed ``argument by argument'' (cf. 4.4.15, p. 99). This, together with the fact that the inclusion $\mathcal{M} \hookrightarrow C^{\wedge}$ is fully faithful, also implies c$'$ $\Longrightarrow$ c.

\msp{page 164}
The implication a$''$ $\Longrightarrow$ c$'$ is seen as in 4.5.6 (p. 119) by applying lemma 4.5.8 1° and 2°, to the case of $C^{\wedge} = \underline{\mathrm{Hom}}(C^{\circ}, \mathcal{E})$, where $\mathcal{E} =$ (Ens). The reverse conditional implication c$'$ $\Longrightarrow$ a$''$, for $\pi \geqslant \pi_1$, is seen as in loc. cit., by applying this lemma 4.5.8. So it still remains to establish the three implications
\[
\left\{
\begin{array}{ccll}
\text{a} & \Longrightarrow & \text{b} & \\[4pt]
\text{a}'' & \Longrightarrow & \text{a} & \text{if } \pi \geqslant \pi_1 \\[4pt]
\text{c} & \Longrightarrow & \text{b} & \text{if } \pi = \pi^{\pi_0}.
\end{array}
\right. \tag*{$(*)$}
\]
But a $\Longrightarrow$ b is tautological, on taking $I = C_0/F$, since we know that $C_0/F$ is large relative to $\pi_0$ and that $F = \varinjlim\limits_{i \in I} x_i$, practically by definition of $\mathcal{M} = \underbrace{\Ind_{\pi}(C)}_{\simeq \Ind_{\pi} C_0} \subset C^{\wedge}$. So it remains to establish the two conditional implications of $(*)$. The implication a$''$ $\Longrightarrow$ a for $\pi \geqslant \pi_1$ is established as in 4.5.2 (cf. page 109), it is essentially sorital. It remains to establish
\[
\text{c} \;\Longrightarrow\; \text{b} \qquad \text{if} \;\; \pi \;=\; \pi^{\pi_0} \;. \tag*{$(*)$}
\]
We will need the auxiliary result:

\msp{page 165}
\textbf{Lemma 4.7.3.} \textit{Let $\pi \geqslant \pi_0$ be such that $\pi^{\pi_0} = \pi$ (thus in fact $\pi > \pi_0$). For every object $F$ of $\mathcal{M}$, we have}
\[
F \;=\; \varinjlim_{i \in I} x_i \qquad \text{(\footnotemark[67])}\;. \tag*{(4.7.3.1)}
\]
\footnotetext[67]{the $x_i$ in $\Filt^{\pi}$, $I$ an ordered set large relative to $\pi$.}
\textit{If in addition $\pi' > \pi$ and $F \in \Ob \Filt^{\pi'}$, then we can take $I$ such that}
\[
\card I \;\leqslant\; \pi'^{\pi} \;. \tag*{(4.7.3.2)}
\]

Indeed, we have
\[
F \;=\; \varinjlim_{j \in J} x_j \qquad \text{with } J = C_0/F \qquad \text{(\footnotemark[68])}\;,
\]
\footnotetext[68]{where $J$ can also be replaced by a category cofinal in [?] $C_0/F$ large relative to $\pi_0$.}
with $J$ large relative to $\pi_0$. Then let
\[
I \;\;=\;\; \parbox[t]{0.62\textwidth}{set of subcategories $J'$ of $J$ \textit{large relative to} $\pi_0$, and such that $\card J' \leqslant \pi$.} \tag*{(4.7.3.3)}
\]

\printedpage{55}
% head: A. Grothendieck : Dérivateurs, Ch. XVIII, Édition des Universités de Montpellier II et Paris VII

I say that $I$ (ordered by inclusion) \textit{is large relative to} $\pi$ (\footnote[69]{cf. proof 4.7.5, page 168-171.}). This follows at once from lemma 4.5.5, p. 113, (stated at first, erroneously, assuming only $\pi > \pi_0$ and not $\pi^{\pi_0} = \pi$). Let us then put, for every $J' \in I$,
\[
x_{J'} \;=\; \varinjlim_{j \in J'} x_j \;,
\]
we thus have
\[
x_{J'} \;\in\; \Ob \Filt^{\pi}
\]
by definition of $I$ and of $\Filt^{\pi}$. On the other hand,
\[
J' \;\longmapsto\; x_{J'} \;:\; I \;\longrightarrow\; \Filt^{\pi}
\]
\msp{page 166}
is a functor, hence
\[
\varinjlim_{J' \in I} x_{J'} \;\simeq\; \varinjlim_{j \in J} x_j \;\simeq\; x \;, % [sic] x -> F, the colimit is the object F of (4.7.3.1)
\]
whence the first assertion of 4.7.3. There remains the bound on cardinals for $F \in \Ob \Filt^{\pi}$. In this case, replacing $C_0/F$ if need be by a cofinal category large relative to $\pi_0$, $\card J \leqslant \pi'$. We then see at once that $\card I \leqslant \pi'^{\pi_0}$, which is the announced bound.

We can now prove the implication $(*)$, page 164, i.e.
\[
\mathcal{M}_{\pi} \;\subset\; \Filt^{\pi} \qquad (\text{for } \pi = \pi^{\pi_0}).
\]
So let $F$ be in $\mathcal{M}_{\pi}$, i.e. $\pi$-accessible, let us write $F$ in the form (4.7.3.2), with $I$ large relative to $\pi$. By definition of $\pi$-accessibility, applied to $F \xrightarrow{\mathrm{id}} F \simeq \varinjlim F_i$, we see that $F$ is a direct factor of an $F_i$, hence of an object of $\Filt^{\pi}$. Not knowing that $\Filt^{\pi}$ is stable under direct factors, we cannot conclude from this that $F$ is in $\Filt^{\pi}$ (so I cannot prove c $\Longrightarrow$ b), but only that $F$ is in $\Filt'^{\pi}$, which, for its part, is stable under $\varinjlim$ large relative to $\pi_0$, of given cardinality $\pi$ $> \pi_0$, hence also \msp{page 167} stable under direct factors. This completes the proof of 4.7.1, and at the same time we get:

% marginal vertical bar beside Corollaire 4.7.4 and its statement
\textbf{Corollary 4.7.4.} \textit{If $\pi \geqslant \pi_0$ is such that $\pi^{\pi_0} = \pi$, then we have}
\[
\Filt'^{\pi} \;=\; (\Filt^{\pi})^{\natural} \;, \tag*{(4.7.4.1)}
\]
\textit{i.e.} $\Filt'^{\pi}$ \textit{consists of the objects $F$ of $\mathcal{M}$ which are direct factors of an object of} $\Filt^{\pi}$.

Indeed, since $\Filt^{\pi} \subset \Filt'^{\pi}$, and since $(\Filt^{\pi})^{\natural} \subset \Filt'^{\pi}$, and since $\Filt'^{\pi}$ is stable under direct factors, we have $(\Filt^{\pi})^{\natural} \subset \Filt'^{\pi}$ for every $\pi \geqslant \pi_0$. But conversely, since $\Filt'^{\pi} \subset \mathcal{M}_{\pi}$, the preceding argument shows that for $\pi = \pi^{\pi_0}$, we have $\Filt'^{\pi} \subset \mathcal{M}_{\pi} \subset (\Filt^{\pi})^{\natural}$, which proves 4.7.3.

\textbf{Remark 4.7.4.1} \texttt{[twice 4.7.4.1]}. Of course, by 4.7.2, 1°, as soon as $\pi \geqslant \pi_1$, we even have $\Filt^{\pi} = \Filt'^{\pi}$ (without having to assume $\pi = \pi^{\pi_0}$), hence \textit{if} $\pi > \pi_0$ (which will be the case if $\pi_1 > \pi_0$, i.e. $\pi_1 \neq \pi_0$). $\Filt'^{\pi}$ being stable under direct factors, $\Filt^{\pi}$ is so too, (4.7.4.1) is satisfied trivially. But if $\pi = \pi_1 = \pi_0$, we have $\Filt^{\pi} = \Filt^{\pi_0} = \Filt'^{\pi_0} =$ subcategory of functors \textit{representable} in $C$ (equivalent to $C$), and the latter is equal to $C^{\natural}$ in $\mathcal{M}$ only if $C$ is Karoubian, cf. \ldots{} below. % [sic] la cas -> le cas, gender of "cas"

\printedpage{56}
% head: A. Grothendieck : Dérivateurs, Ch. XVIII, Édition des Universités de Montpellier II et Paris VII

\msp{page 168}
\textbf{4.7.5.} For the proof of 4.7.1 to be complete, I would still have to give a proof of 4.5.5, p. 113, which is the key technical lemma for establishing the

\textbf{Lemma 4.7.5.1} (used in the proof of 4.7.3). \textit{Let $\pi_0$, $\pi$ be two infinite cardinals, with $\pi \geqslant \pi_0$ and $\pi^{\pi_0} = \pi$ (hence $\pi > \pi_0$). Let $J$ be a category large relative to $\pi_0$. Then the set $I$, ordered by inclusion, of subcategories $J'$ of $J$ which are large relative to $\pi_0$, and such that} $\card J' \leqslant \pi$, \textit{is large relative to $\pi$. (Moreover, if} $\card J \leqslant \pi'$, \textit{we have} $\card I \leqslant \pi'^{\pi}$\textit{.)}

So we must show the

\textbf{Corollary 4.7.5.2} (\footnote[70]{Compare with the proof of 4.6.1.2.}) (this is lemma 4.5.5 not yet proved). \textit{Under the conditions of 4.7.5, let $K \subset \mathrm{Fl} J$ be such that} $\card K \leqslant \pi$. \textit{Then there exists $J' \in I$ such that $K \subset J'$.}

We know, $\pi$ being infinite, that the category $K'$ generated by $K$ is also of cardinality $\leqslant \pi$, so we can suppose that $K$ ``is'' a subcategory of $J$. We are going to construct a $J' \in I$ containing $K$, by constructing by transfinite recursion an increasing sequence $K_{\alpha}$ of subcategories of $J$, with $\card K_{\alpha} \leqslant \pi$, indexed by the ordinals $\alpha$ such that \fbox{$\card \alpha \leqslant \pi$}:

\msp{page 169}
\begin{enumerate}
\item[1°)] $K_0 = K$.
\item[2°)] $K_{\alpha} = \bigcup_{\beta < \alpha} K_{\beta}$ ($= \varinjlim K_{\beta}$) if $\alpha$ is a limit ordinal (hence $\card K_{\alpha} \leqslant \pi$, since the $\card K_{\beta} \leqslant \pi$).
\item[3°)] $K_{\alpha+1}$ obtained from $K_{\alpha}$ in the following way. If $K_{\alpha}$ is large relative to $\pi_0$, we take $K_{\alpha+1} = K_{\alpha}$ (hence, in view of 2°, we will have $K_{\alpha'} = K_{\alpha}$ for $\alpha' \geqslant \alpha$ -- the transfinite construction stops \ldots). Otherwise, $K_{\alpha+1}$ is obtained as follows:
\begin{enumerate}
\item[a)] for every subset $A$ of cardinality $\leqslant \pi_0$ of $\Ob K_{\alpha}$, we choose an upper bound $x_A$ in $J$ (we feel like taking it in $K_{\alpha}$, if there is an upper bound in $K_{\alpha}$)\,;
\item[b)] for every pair $x$, $y$ of objects of $K_{\alpha}$, and every subset of cardinality $\leqslant \pi_0$ of $\mathrm{Hom}_{K_{\alpha}}(x,y)$, we choose an arrow $u_B : y \longrightarrow z_B$ in $J$, which equalizes the $v_{\beta} \in B$ (and we take this arrow in $K_{\alpha}$, when possible)\,;
\item[c)] we take for $K_{\alpha+1}$ the smallest subcategory containing the $\mathrm{id}_{x_A}$ and the $u_B$.
\end{enumerate}
\end{enumerate}

We must check that
\[
\card K_{\alpha+1} \;\leqslant\; \pi \;.
\]
The number of possible choices for $A \in \mathfrak{P}_{\leqslant \pi_0}(\Ob J)$ \msp{page 170} is $\leqslant \pi^{\pi_0} = \pi$, the number of possible choices for $B \in \mathfrak{P}_{\pi_0}(\mathrm{Fl}\, J)$ is also $\leqslant \pi^{\pi_0} = \pi$, so the number of arrows of the form $\mathrm{id}_{x_A}$ or $u_B$ is $\leqslant \pi + \pi = \pi$. But then the subcategory $K_{\alpha+1}$ of $J$ generated by this set of arrows is of cardinality $\leqslant 1$. Then let % [sic] \leqslant 1 -> \leqslant \pi, the statement to be verified is card K_{\alpha+1} \leqslant \pi
\printedpage{57}
% head: A. Grothendieck : Dérivateurs, Ch. XVIII, Édition des Universités de Montpellier II et Paris VII
$\alpha_0$ be the least ordinal such that
\[
\card \alpha_0 \;>\; \pi_0 \;,
\]
we thus have
\[
\card \alpha_0 \;=\; \text{successor cardinal of } \pi_0 \leqslant \pi \;,
\]
since $\pi^{\pi_0} = \pi$, whence $\pi > \pi_0$ (for $\pi^{\pi_0} = \pi_0^{\pi_0} \geqslant 2^{\pi_0} > \pi_0$). Thus, $\alpha_0$ is among the admissible ordinals $\alpha$ (i.e. such that $\card \alpha \leqslant \pi$). I say that $K_{\alpha_0}$ \textit{is large relative to} $\pi_0$ (which will prove 4.7.5.2, on taking $J' = K_{\alpha_0}$). This follows immediately from the fact that
\[
K_{\alpha_0} \;=\; \bigcup_{\alpha < \alpha_0} K_{\alpha} \qquad (\text{\textit{increasing} union}),
\]
% double marginal vertical bar beside the following two lines
and that the set of $\alpha < \alpha_0$ is \textit{large relative to} $\pi_0$ (\footnote[71]{spell this out in a

\textbf{Lemma.} \textit{Let $\pi_0$ be a cardinal, $\alpha_0 = \alpha_{\pi'_0}$ ($\pi'_0$ the successor of $\pi_0$) the least ordinal such that} $\card \alpha_0 > \pi_0$, \textit{then $I_{\alpha_0} = \{\alpha \mid \alpha < \alpha_0\}$ is large relative to $\pi_0$.}}). So every set of objects or arrows of $K_{\alpha_0}$, of cardinality $\leqslant \pi_0$, is contained in some $K_{\alpha}$ with $\alpha < \alpha_0$.

\msp{page 171}
This implies, if it is a set of objects, that there is a common upper bound in $K_{\alpha+1}$, a fortiori in $K_{\alpha_0}$. If it is a subset of a $\mathrm{Hom}_{K_{\alpha}}(x,y)$, that there is an equalizer $y \longrightarrow z$ in $K_{\alpha+1}$, a fortiori in $K_{\alpha_0}$. Thus the criterion $\mathrm{PF}_{\pi}$ 1,2 of 4.4.4 (page 89) is satisfied by $K_{\alpha_0}$, which completes the proof of 4.7.5.2.

\textbf{Corollary 4.7.6} (of 4.7.1). \textit{For every cardinal $\pi \geqslant \pi_0$, we have}
\[
\underbrace{\Filt'^{\pi}}_{\overset{\textit{if } \pi \geqslant \pi_1}{=}\ \Filt^{\pi}} \;\subset\; \mathcal{M}_{\pi} \;\subset\; \underbrace{\Filt^{\pi^{\pi_0}}}_{\overset{\textit{if } \pi^{\pi_0} \geqslant \pi_1}{=}\ \Filt^{\pi^{\pi_0}}} \qquad (= (\Filt^{\pi^{\pi_0}})^{\natural} \textit{ by cor. 4.7.4})\,, \tag*{(4.7.6.1)} % ?? uncertain: tiny condition under the second brace -- "si \pi^{\pi_0} \geqslant \pi_1" or "si \pi^{\pi}_0 \geqslant \pi_1"
\]
\textit{therefore if $\pi \geqslant \pi_1$, we also have}
\[
\Filt^{\pi} \;\subset\; \mathcal{M}_{\pi} \;\subset\; \Filt^{\pi^{\pi_0}} \qquad \textit{(if $\pi^{\pi_0} \geqslant \pi_1$).} \tag*{(4.7.6.2)}
\]

Indeed, $\Filt'^{\pi} \subset \mathcal{M}_{\pi}$ is none other than the implication b$'$ $\Longrightarrow$ c of 4.7.1, on the other hand, we have $\pi \leqslant \pi^{\pi_0}$, whence
\[
\mathcal{M}_{\pi} \;\subset\; \mathcal{M}_{\pi^{\pi_0}} \;=\; \Filt'^{\pi^{\pi_0}} \;,
\]
where the last equality comes from $\mathcal{M}_{\pi'} = \Filt'^{\pi'}$ if $\pi'^{\pi_0} = \pi'$, i.e. from (4.7.2.3) (page 163), putting $\pi' = \pi^{\pi_0}$, taking into account that $\pi'^{\pi_0} = (\pi^{\pi_0})^{\pi_0} = \pi^{\pi_0 \cdot \pi_0} = \pi^{\pi_0} = \pi$.

\msp{page 172}
\textbf{Proposition 4.7.7.} \textit{Consider $C$, $C_0$, $\mathcal{M} = \Ind_{\pi_0}(C)$, $\pi_1 = \sup(\pi_0, \card \mathrm{Fl}\, C_0)$ as before. For every cardinal $\pi \geqslant \pi_0$, let $\mathcal{F}^{\pi}$ be the strictly full subcategory} $\Filt'^{\pi}$ \textit{(resp. $(C^{\wedge})_{\pi} \cap \mathcal{M}$, resp. $\mathcal{M}_{\pi}$) in $\mathcal{M}$. Then the increasing filtration $(\mathcal{F}^{\pi})_{\pi \geqslant \pi_0}$ of $\mathcal{M}$ by the subcategories $\mathcal{F}^{\pi}$ satisfies the following conditions.}
\begin{enumerate}
\item[\textit{a)}] \textit{For every $\pi \geqslant \pi_0$, $\mathcal{F}^{\pi}$ is equivalent to a small category.}
\item[\textit{b)}] \textit{$\mathcal{M}$ satisfies $\mathrm{L}_{\pi_0}$ (i.e. is stable under filtering $\varinjlim$ large relative to $\pi_0$), and for every $\pi \geqslant \pi_0$, $\mathcal{F}^{\pi}$ is stable under filtering limits large relative to $\pi_0$ and of cardinality $\leqslant \pi_0$ [?].}
\printedpage{58}
% head: A. Grothendieck : Dérivateurs, Ch. XVIII, Édition des Universités de Montpellier II et Paris VII
\item[\textit{c)}] \textit{If $\pi \geqslant \pi_0$ is} such that \fbox{$\pi^{\pi_0} = \pi$,} \textit{then for every $F$ in $\mathcal{M}$, we have an isomorphism}
\[
F \;=\; \varinjlim_{i \in I} x_i \;,
\]
\textit{the $x_i$ in $\mathcal{F}^{\pi}$, and $I$ an ordered category large relative to $\pi_0$. If in addition $\pi' > \pi$ and $F \in \Ob F^{\pi'}$, we can take $I$ such that} % [sic] F^{\pi'} (italic F) -> \mathcal{F}^{\pi'}, the filtration is denoted \mathcal{F}^{\pi}
\[
\card I \;\leqslant\; \pi'^{\pi} \qquad \underline{\textit{if}}\ \pi'^{\pi} \geqslant \pi_1. \tag*{(4.7.7.1)}
\]
\item[\textit{d)}] \textit{If $\pi \geqslant \pi_0$, we have $\mathcal{F}^{\pi} \subset \mathcal{M}_{\pi}$, i.e. the objects of $\mathcal{F}^{\pi}$ are $\pi$-accessible.}
\end{enumerate}
\textbf{N.B.} These are exactly, apart from d), the conditions a), b), c) for cardinal filtrations of SGA 4 I 9.7.12, \msp{page 173} except only that in c), we restrict ourselves here to the ordinals $\pi$ such that $\pi^{\pi_0} = \pi$ (\footnote[72]{and \textit{in addition}, in the bound (4.7.7.1), to the cardinals $\pi' > \pi$ such that we also have $\pi'^{\pi} \geqslant \pi_1$ (which is automatic if we have $\pi \geqslant \pi_1$, which is the case, given that $\pi = \pi^{\pi_0} \geqslant 2^{\pi_0}$, if we have $2^{\pi_0} \geqslant \pi_1$.)}). The additional condition d) is an ``accessibility condition'', cf. below.

\textsc{Proof of} 4.7.7. a) Since we have
\[
\Filt'^{\pi} \;\subset\; (C^{\wedge})_{\pi} \cap \mathcal{M} \;\subset\; \mathcal{M}_{\pi} \;,
\]
we are reduced to proving that $\mathcal{M}_{\pi}$ is equivalent to a small category. We are reduced to the case $\pi \geqslant \pi_1$, $\pi = \pi^{\pi_0}$, in which case we have $\mathcal{M}_{\pi} = \Filt^{\pi}$. In fact, we find:

\textbf{Corollary 4.7.7.1} \texttt{[twice 4.7.7.1]} (\footnote[73]{Compare with 4.7.19.}). \textit{If $\pi \geqslant \pi_1$, we have}
\[
\left\{
\begin{array}{rcll}
\cardred \Filt^{\pi} & \leqslant & 2^{\pi} & \\[4pt]
\cardred \Filt'^{\pi} & \leqslant & \cardred \mathcal{M}_{\pi} \;\leqslant\; 2^{(\pi^{\pi_0})} & \qquad \textit{($= 2^{\pi}$ if $\pi^{\pi_0} = \pi$).}
\end{array}
\right.
\]

The second bound follows from the first and from
\[
\mathcal{F}^{\pi} \;\subset\; \mathcal{M}_{\pi} \;\subset\; \Filt^{\pi^{\pi_0}} \qquad (\text{for } \pi^{\pi_0} \geqslant \pi_1) \qquad \text{(\footnotemark[74]).}
\]
\footnotetext[74]{Cf. (4.7.6.2).}
For the first, we use condition a$''$) of 4.7.1 (which characterizes $\Filt^{\pi}$ for $\pi \geqslant \pi_1$). I pass over the details of the verification.

\msp{page 174}
Let us prove 4.7.2 b). The condition $\mathrm{L}_{\pi_0}$ is known, there remains the stability of $\mathcal{F}^{\pi}$. In the case of $\Filt'^{\pi}$, it is established by definition. In the cases $C^{\wedge}_{\pi} \cap \mathcal{M}$ and $\mathcal{M}_{\pi}$, it follows from the stability of the notion of $\pi$-accessibility under $\varinjlim$ of cardinality $\leqslant \pi$, either in $C^{\wedge}$ or in $\mathcal{M}$ (taking into account that $C^{\wedge}$ and $\mathcal{M}$ have a generating family $C$ formed of $\pi_0$-accessible objects, hence $\pi$-accessible).

Let us prove c): this is essentially lemma 4.7.3, taking into account that we have
\[
\Filt^{\pi} \;\subset\; \mathcal{F}^{\pi} \;,
\]